\documentclass[a4paper,12pt,fleqn]{scrartcl}%fleqn pour aligner les equations à gauche
\usepackage{../../Styles/Cours}

\newcommand{\chnum}{Chapitre 4}
\newcommand{\chtitle}{Structure de l'atome}
\newcommand{\dtype}{Activité}

\begin{document}
	\maketitle

\noindent \begin{minipage}{\linewidth}
	\begin{documentboxenv}{La structure de l'atome}
	« Lorsque j'entrais au laboratoire dirigé par Joliot au Collège de France, la connaissance que j'avais de la structure de la matière ne devait guère dépasser celle acquise par un lycéen de 1993 abonné à de bonnes revues de vulgarisation. Je les résume rapidement : la matière est composée d'atomes, eux-mêmes constitués de noyaux entourés d'un cortège d'électrons. Les noyaux portent une charge électrique positive qui est de même valeur et de signe opposé à la charge des électrons qui gravitent autour du noyau. La masse d'un atome est concentrée dans le noyau. (...)\\
	Le noyau de l'hydrogène ou proton, porte une charge électrique positive. Celui-ci a un compagnon, le neutron, qui est neutre électriquement et a sensiblement la même masse. Tous deux s'associent de façon très compacte pour constituer les noyaux qui sont au cœur des atomes peuplant notre univers. Ils s'entourent d'un cortège d'électrons dont la charge compense exactement celle des protons. En effet, la matière est neutre, sinon elle exploserait en raison de la répulsion qu'exercent l'une sur l'autre des charges de même signe, positif ou négatif. Il faut avoir en tête l'échelle des dimensions. Le diamètre d'un atome est voisin d'un centième de millionième de centimètre. Celui d'un noyau est cent mille fois plus petit. On voit donc que presque toute la masse d'un atome est concentrée en un noyau central et que, loin sur la périphérie, se trouve un cortège qui est fait de particules de charge électrique négative, les électrons. C'est ce cortège seul qui gouverne le contact des atomes entre eux et donc tous les phénomènes perceptibles de notre vie quotidienne, tandis que les noyaux, tapis au c\oe ur des atomes, en constituent la masse. »
	\begin{flushright}
		Georges CHARPAK, Extrait de \emph{La vie à fil tendu}
	\end{flushright}
	\end{documentboxenv}
\end{minipage}

\vspace{1em}
\noindent\begin{minipage}[t]{.55\linewidth}
	\begin{documentboxenv}{La notation symbolique}
	Le noyau d'un atome est représenté de manière symbolique comme suit : $^A_ZX$
	\begin{itemize*}
		\item $\mathbf{X}$ est le symbole chimique de l'atome
		\item $\mathbf{A}$ est le nombre de nucléons du noyau
		\item $\mathbf{Z}$ est le numéro atomique correspondant au nombre de protons dans le noyau
	\end{itemize*}
Le nombre de neutrons $\mathbf{N}$ est calculé par : $\mathbf{N = A - Z}$.\\
\textcolor{gray!80}{Exemple : la notation symbolique du noyau d'azote est $^{14}_7N$ constitué de 7 protons et 7 neutrons.}\\
Deux noyaux peuvent avoir le même nombre de protons mais un nombre de neutrons différents, ils sont appelés des \textbf{isotopes.}\\
\textcolor{gray!80}{Exemple : il existe plusieurs isotopes du carbone dont les notations symboliques sont : $^{12}_6C$, $^{13}_6C$ et $^{14}_6C$.}
\end{documentboxenv}\end{minipage}
\hspace{.02\linewidth}
\begin{minipage}[t]{.43\linewidth}
	\begin{documentboxenv}{Données supplémentaires }
	Masse d'un électron : \SI{9,1094e-31}{kg}\\
	Masse d'un neutron : \SI{1,6749e-27}{kg}\\
	Masse d'un proton : \SI{1,6726e-27}{kg}
	\end{documentboxenv}
\end{minipage}\\
\newpage
\textbf{Questions :}
\begin{enumerate}
	\item Compléter le tableau suivant :
\end{enumerate}
\begin{center}
	\begin{tabular}{|>{\centering}m{4cm}|m{4.5cm}|m{4.5cm}|m{4.5cm}|}
	\hline
	\textbf{Particules citées dans le texte}&&&\\
	\hline
	\textbf{Où les trouve-t-on (noyau, atome) ?}&&&\\
	\hline
	\textbf{Quelle est leur charge électrique ?}&&&\\
	\hline
	\textbf{Quelle est leur masse ?}&&&\\
	\hline 
\end{tabular}
\end{center}
\begin{enumerate}[resume]
	\item Comparer la masse d'un nucléon à celle d'un électron. Quelle peut-on en conclure quant à la répartition de la masse dans les atomes ?
	\item De quoi est constitué l'atome d'hydrogène ? Donner la notation symbolique du noyau d'hydrogène.
	\item Quelle est la taille approximative d'un noyau (en mètre) ? Quelle est celle d'un atome (en mètre ) ?
	\item Donner la constitution des noyaux suivants et noter vos remarques concernant les noyaux de vanadium (\ce{V}).
\end{enumerate}
\begin{center}
	\begin{tabular}{|>{\centering}m{4cm}|>{\centering}m{3cm}|>{\centering}m{3cm}|>{\centering}m{3cm}|>{\centering\arraybackslash}m{3cm}|}
	\hline
	\textbf{Noyau} & $^{12}_6C$ & $^{24}_{12}Mg$ &$^{50}_{23}V$ &$^{55}_{23}V$ \\
	\hline
	\textbf{Nombre de protons}&&&&\\
	\hline
	\textbf{Nombre de nucléons}&&&&\\
	\hline
	\textbf{Nombre de neutrons}&&&&\\
	\hline
\end{tabular}
\end{center}


\begin{enumerate}[resume]
	\item Donner la composition d'un atome d'oxygène $^{16}_8O$. Calculer la masse de son noyau et de l'atome puis conclure.
	\item \underline{Pour les plus rapides :} \emph{on dit que "comparer la taille d'un noyau d'atome est celle d'un atome revient à comparer la taille d'une pointe de stylo (0,5 mm) au milieu d'un terrain de football (100 m de long sur 45 m de large)". Est-ce que cette comparaison est fondée ? (utiliser un schéma)}
\end{enumerate}
\end{document}